\section{Formal Results for the HIF Score Under Stated Assumptions}\label{app:formal_results}

We provide a formal analysis of Hybrid Integrity Framework (HIF) score properties under the assumptions stated below. The results should be interpreted as assumption-scoped statements rather than unconditional guarantees.

\textbf{Theorem 1 (Intersection Soundness).} The HIF score $H(\mathbf{x})$ defines a membership certificate for the intersection of semantic manifolds $\mathcal{M}_{total} = \bigcap_{l} \mathcal{M}_l$. Under the idealization that each component validity satisfies $\mathcal{C}_l(\mathbf{x}) = 1$ iff $\mathbf{x} \in \mathcal{M}_l$ and $\mathcal{C}_l(\mathbf{x}) \in [0, 1)$ otherwise, $H(\mathbf{x}) = \left( \prod_l \mathcal{C}_l(\mathbf{x}) \right)^{1/L}$ attains its maximum value $1$ exactly on $\mathcal{M}_{total}$:
\begin{equation}
    \mathbb{1}[H(\mathbf{x}) = 1] = \prod_{l} \mathbb{1}[\mathbf{x} \in \mathcal{M}_l]
\end{equation}

\textbf{Proof.} \emph{Completeness:} if $\mathbf{x} \in \mathcal{M}_{total}$ then $\mathcal{C}_l(\mathbf{x}) = 1$ for all $l$, so $H(\mathbf{x}) = 1$. \emph{Soundness (non-compensatory property):} if $\mathbf{x} \notin \mathcal{M}_{total}$, there exists $l^*$ with $\mathcal{C}_{l^*}(\mathbf{x}) < 1$, and since all validities lie in $[0,1]$, the geometric mean satisfies $H(\mathbf{x}) < 1$. \emph{Strict convergence:} an ``Absolute Violation'' in any layer ($\mathcal{P}_{l^*} \to 1$) drives $H(\mathbf{x}) \to 0$ regardless of other layers. $H(\mathbf{x})$ thus behaves as a conservative diagnostic for joint satisfaction of latent manifold constraints, matching the implementation in Algorithm~\ref{alg:audit}.

\textbf{Theorem 2 (Asymptotic Convergence of the HIF Penalty).} Let $\hat{P}_N(x_h \mid \mathbf{x}_{\setminus h})$ be the empirical conditional probability estimator (the Sentinel $f_h$) trained on a ground-truth dataset $\mathcal{D}_{real}$ of size $N$. Assuming the base estimator is statistically consistent (e.g., consistent non-parametric trees), as $N \to \infty$, $\hat{P}_N$ converges in probability to the true conditional distribution $P^*(x_h \mid \mathbf{x}_{\setminus h})$, and the empirical penalty $\mathcal{P}_h^{(N)}(\mathbf{x})$ converges pointwise in probability to the corresponding population penalty:
\begin{equation}
    \lim_{N \to \infty} P\left( \left| \mathcal{P}_h^{(N)}(\mathbf{x}) - \mathcal{P}_h^*(\mathbf{x}) \right| > \epsilon \right) = 0 \quad \forall \epsilon > 0, \; \mathbf{x} \in \mathcal{M}_{real}
\end{equation}

\textbf{Proof sketch.} \emph{(i)} Because the LSE leverages consistent non-parametric estimators, the empirical conditional probabilities converge pointwise in probability to the true Bayesian posterior $P^*$ (provided the feature space is bounded and leaf nodes grow sub-linearly with $N$); the pointwise statement reflects that sup-norm convergence is not generally guaranteed for forest-based estimators. \emph{(ii)} Defining the confidence floor $\delta_h$ as the empirical 5th-percentile of the sentinel's out-of-bag probabilities yields an asymptotic upper bound of approximately $\alpha = 0.05$ on false positives under calibration assumptions. \emph{(iii)} Records penalized by the converged HIF ($\mathcal{P}^* > 0$) are expected to lie in lower-probability conditional regions of the manifold, supporting a structural-violation interpretation. \emph{(iv)} By the Uniform Law of Large Numbers and the Continuous Mapping Theorem, the empirical Feature Dependency Score $S^{(N)}(h)$ converges in probability to the Bayes optimal classification accuracy $S^*(h)$, implying asymptotically stable hub ranking.

\textbf{Theorem 3 (Type I Error Bound \& Cohort-Purity Approximation).} Let $\mathbf{x} \sim P_{real}$ be a valid record on the ground-truth manifold $\mathcal{M}_{real}$ (including rare tail outliers). Under quantile calibration ($\delta_h = Q_{\alpha}$ and continuous tolerance bandwidth $\gamma_y = Q_{1-\beta}$), the false-rejection probability is upper-bounded by the Union Bound:
\begin{equation}
    P\left( H(\mathbf{x}) < 0.5 \;\middle|\; \mathbf{x} \in \mathcal{M}_{real} \right) \le \sum_{h \in \mathcal{H}} \alpha_h + \sum_{y} \beta_y
\end{equation}
Furthermore, let a generative model produce a noisy mixture $P_{syn} = (1 - \eta) P_{real} + \eta P_{hallucinated}$, where $\eta \in [0, 1]$ is the structural hallucination rate. Rejection filtering with threshold $H \ge 0.5$ yields an approximate retained-cohort purity:
\begin{equation}
    \text{Purity}(H \ge 0.5) = \frac{(1 - \eta)(1 - \alpha)}{(1 - \eta)(1 - \alpha) + \eta \beta_{hallucination}}
\end{equation}
where $\beta_{hallucination}$ is the leakage probability of hallucinated records past the filter. For high hallucination rates ($\eta \approx 0.66$), the pre-filter purity is $(1 - \eta) = 34\%$; with a calibrated leakage rate $\beta_{hallucination} \approx 0.02$ and $\alpha = 0.05$, filtering elevates purity to over $96\%$ in this illustrative model, theoretically accounting for the downstream predictive utility improvements observed.

\section{End-to-End Auditing Procedure}\label{app:calibration}

Algorithm~\ref{alg:calibration} summarizes the complete manifold calibration procedure---hub selection, confidence-floor calibration, NIC regressor training, and rule mining---and Algorithm~\ref{alg:audit} details the end-to-end auditing procedure for scoring a synthetic dataset with the trained HIF components.

\begin{algorithm}[ht]
\caption{HIF Manifold Calibration}
\label{alg:calibration}
\begin{algorithmic}[1]
\Statex \textbf{Input:} Ground-truth data $\mathcal{D}_{real}$, Hub count $K_{hub}$
\Statex \textbf{Output:} Sentinels $\{f_h\}$, Floors $\{\delta_h\}$, NIC Regressors $\{r_y\}$, Thresholds $\{\tau_y\}$
\State $\mathcal{H} \leftarrow \text{Select top } K_{hub} \text{ features by Feature Dependency Score } S(h)$
\For{each hub $h \in \mathcal{H}$}
    \State Train Sentinel $f_h: \mathcal{X}_{\setminus h} \to [0, 1]^{|\mathcal{X}_h|}$ on $\mathcal{D}_{real}$ to predict $x_h$
    \State $\delta_h \leftarrow \max( \text{Percentile}_{5} ( \{ f_h^{\text{OOB}}(\mathbf{x}_h \mid \mathbf{x}_{\setminus h}) \mid \mathbf{x} \in \mathcal{D}_{real} \} ), 0.01)$
\EndFor
\State $\mathcal{Z} \leftarrow \text{SVD}(\text{Scale}(\text{OneHot}(\mathcal{D}_{cat})))$ \Comment{Non-Centered Spectral Manifold}
\For{each continuous feature $y$}
    \State Train Gradient Boosting regressor $r_y: \mathcal{Z} \to \mathbb{R}$ on $\mathcal{D}_{real}$
    \State $\mathcal{E}_y \leftarrow \{ |y - r_y(\mathbf{z})| \mid \mathbf{x} \in \mathcal{D}_{real} \}$ \Comment{Ground-truth residuals}
    \State $\tau_y \leftarrow \max(Q_{0.95}(\mathcal{E}_y), \text{Median}(\mathcal{E}_y) + 2 \cdot \text{MAD}(\mathcal{E}_y))$ \Comment{Robust Threshold}
    \State $\gamma_y \leftarrow \max(Q_{0.98}(\mathcal{E}_y), 2 \cdot \tau_y, 3 \cdot \text{MAD}(\mathcal{E}_y), 0.1)$ \Comment{Dynamic Gamma Scaling}
\EndFor
\State $\mathcal{R} \leftarrow \text{MineImplicationRules}(\mathcal{D}_{real})$ \Comment{Apriori-style, min\_conf=0.95, min\_supp=0.005}
\end{algorithmic}
\end{algorithm}

\begin{algorithm}[ht]
\caption{HIF Synthetic Data Audit}
\label{alg:audit}
\begin{algorithmic}[1]
\Statex \textbf{Input:} Synthetic record $\mathbf{x}$, Trained components from Alg \ref{alg:calibration}
\Statex \textbf{Output:} Record Integrity Score $H(\mathbf{x})$
\State \textbf{Step 1: Logical Sentinel Audit (Categorical)}
\For{each hub $h \in \mathcal{H}$}
    \State $\mathcal{P}_h \leftarrow \text{clip}\left( \frac{\delta_h - f_h(\mathbf{x}_h \mid \mathbf{x}_{\setminus h})}{\delta_h}, 0, 1 \right) $
\EndFor
\State $\mathcal{P}_{cat} \leftarrow 1 - \prod_{h \in \mathcal{H}} (1 - \mathcal{P}_h)$ \Comment{Multiplicative violation aggregation}
\State \textbf{Step 2: Neighbor-Invariant Continuity Audit (Continuous)}
\State $\mathbf{z} \leftarrow \phi(\mathbf{x}_{cat})$ \Comment{Apply learned standardized projection}
\For{each continuous feature $y$}
    \State $\rho_y \leftarrow |y - r_y(\mathbf{z})|$
    \State $\mathcal{P}_y \leftarrow \text{clip}\left( \frac{\rho_y - \tau_y}{\gamma_y}, 0, 1 \right)$
\EndFor
\State $\mathcal{P}_{cont} \leftarrow \max_y \mathcal{P}_y$
\State \textbf{Step 3: Integrity Aggregation}
\State $\mathcal{P}_{rule} \leftarrow 1 \text{ if } \mathbf{x} \text{ violates any } r \in \mathcal{R} \text{ else } 0$
\State $H(\mathbf{x}) \leftarrow \left( \prod_{l} (1 - \mathcal{P}_l) \right)^{1/L} \in [0, 1]$ \Comment{Geometric mean of validities}
\end{algorithmic}
\end{algorithm}

\section{Logical Sentinel Ensemble (LSE) Implementation Details}\label{app:lse_details}

On the U.S. Census ACS dataset, the Feature Dependency Score $S(h)$ identified the following top-5 hubs: \texttt{household\_income}, \texttt{housing\_cost}, \texttt{poverty\_status}, \texttt{education}, and \texttt{tenure}. The confidence floor $\delta_h$ is the 5th-percentile of the \emph{out-of-bag} probability mass that a sentinel assigns to the correct category in the ground-truth data (safety floor 0.01); out-of-bag predictions yield out-of-sample floors without a held-out split, reflecting generalization rather than memorized fit, and avoid the instability of 1st-percentile thresholds in sparse data \cite{karabulut2025neurosymbolic, long2025evaluating}. Hub selection is deterministic: sentinel discovery uses a fixed random state with non-shuffled cross-validation folds, so the hub set depends only on the ground-truth data and is verified identical across all $N=10$ held-out seeds for every benchmark dataset (e.g., \texttt{household\_income} through \texttt{tenure} on Census ACS; \texttt{quantity} through \texttt{item\_tax} on Online Purchases).

\section{Extended Results and Sensitivity Analysis}

\subsection{Hyperparameter Sensitivity}\label{app:hyperparameters}

We evaluated HIF robustness across its three key hyperparameters on the Census ACS dataset (2,000 records with 5\% injected violations; $N=5$ seeds, mean $\pm$ SEM). The continuous HIF score is the primary metric; the violation rate derives from applying the violation threshold to row-level penalties. The HIF score is stable across hub counts (varying by $<1\%$ from $0.979 \pm 0.0003$ at $K=1$ to $0.987 \pm 0.0007$ at $K\geq 10$), since geometric-mean aggregation prevents added sentinels from diluting the score; the hub set saturates at $K=10$ because Census ACS contains only 10 features. Scores are likewise stable for confidence-floor percentiles 1--5 ($0.9848 \pm 0.0002$ at percentile 1 vs.\ $0.9841 \pm 0.0002$ at percentile 5; minor $0.3\%$ drop to $0.9815 \pm 0.0002$ at the 10th), making the 5th-percentile default a good sensitivity--robustness balance. Finally, the continuous score is invariant to the binary threshold ($0.9841 \pm 0.0002$ for thresholds 0.3, 0.5, 0.7), while the violation rate scales monotonically ($0.81\% \to 0.37\% \to 0.21\%$), letting practitioners tune the operating point without retraining. The same threshold sweep extends to the downstream-utility protocol (Section~\ref{sec:utility}): the integrity-filtered F1 gain on Census ACS CTGAN is significant at every operating point and increases monotonically with strictness (Table~\ref{tab:threshold_utility}).

\begin{table}[htbp]
\centering
\caption{Downstream-Utility Threshold Sensitivity (Census ACS, CTGAN, $N=10$ seeds, mean). Integrity-filtered TSTR F1 and retention at each frontier threshold $H$; the unfiltered cohort achieves F1 $= 0.427$. The recovery effect is significant at every operating point and grows with filtering strictness.}
\label{tab:threshold_utility}
\footnotesize
\setlength{\tabcolsep}{3pt}
\begin{tabular}{@{}lcccccc@{}}
\toprule
Threshold $H$ & Retained & F1 & $\Delta$F1 & Paired $t$ $p$ & Wilcoxon $p$ & 95\% CI for $\Delta$F1 \\
\midrule
0.3 (strict) & 30.9\% & 0.723 & $+0.295$ & $<0.001$ & 0.002 & $[+0.215, +0.375]$ \\
0.5 (default) & 65.7\% & 0.531 & $+0.104$ & 0.002 & 0.002 & $[+0.049, +0.160]$ \\
0.7 (loose) & 92.6\% & 0.451 & $+0.023$ & 0.020 & 0.037 & $[+0.005, +0.042]$ \\
\bottomrule
\end{tabular}
\end{table}

\subsection{Sample-Complexity Floor at HIF's Operating Configuration}\label{app:sample_complexity}

These bounds are computable at HIF's operating configuration. The identity/uniformity-testing bound $m \gtrsim \sqrt{S}/\varepsilon^2$ \cite{paninski2008,valiant2014instance} inverted at the per-hub, per-category sample sizes HIF actually conditions on yields per-cell detectability floors $\varepsilon_{\text{cell}} = (S_{\text{cell}} / n_{\text{cell}}^2)^{1/4}$, where $n_{\text{cell}}$ is the number of real rows with hub value $c$ and $S_{\text{cell}}$ is the distinct joint configurations of the non-hub features in that cell. Table~\ref{tab:sample_complexity} reports these values for each benchmark dataset using the auditor's default configuration (5 hubs, 10 quantile bins for continuous features). The cohort-level analogue $\varepsilon_{\text{joint}} = (S_{\text{joint}} / n^2)^{1/4}$ with $S_{\text{joint}}$ distinct full-row configurations and $n=2000$ is also shown.

\begin{table}[htbp]
\centering
\caption{Sample-Complexity Floor: per-cell detectability $\varepsilon_{\text{cell}}$ inverted from the identity-testing bound at HIF's actual hub-conditioned sample sizes (5 hubs per dataset, 10 quantile bins). The worst cell (smallest $n_{\text{cell}}$) drives the floor; aggregate $\varepsilon_{\text{joint}}$ assumes a global uniformity test on the full joint.}
\label{tab:sample_complexity}
\footnotesize
\setlength{\tabcolsep}{2.5pt}
\begin{tabular}{@{}lrrrrrrr@{}}
\toprule
Dataset & Rows & Med. $n_{\text{cell}}$ & Min $n_{\text{cell}}$ & Med. $\varepsilon_{\text{cell}}$ & Worst $\varepsilon_{\text{cell}}$ & $\varepsilon_{\text{joint}}$ \\
\midrule
Census ACS & 2000 & 200 & 198 & 0.266 & 0.267 & 0.150 \\
Adult & 2000 & 5.5 & 1 & 0.654 & 1.000 & 0.149 \\
Credit & 2000 & 12 & 1 & 0.537 & 1.000 & 0.150 \\
Online Purchases & 664 & 66 & 1 & 0.325 & 1.000 & 0.193 \\
Supermarket Sales & 1000 & 100 & 100 & 0.316 & 0.316 & 0.178 \\
\bottomrule
\end{tabular}
\end{table}

The worst-case floors ($\varepsilon \approx 1.0$) occur where long-tail hub categories leave cells with $n_{\text{cell}} \le 12$ (Adult \texttt{native\_country}, Credit \texttt{pay\_*}, Online Purchases \texttt{quantity}). Census ACS, whose hubs (\texttt{household\_income}, \texttt{housing\_cost}, \texttt{poverty\_status}, \texttt{education}, \texttt{tenure}) all have balanced category counts ($n_{\text{cell}} \ge 198$), achieves $\varepsilon \approx 0.27$; Supermarket Sales' hubs are likewise balanced ($n_{\text{cell}} = 100$, $\varepsilon \approx 0.32$). The cohort-level analogue $\varepsilon_{\text{joint}} \approx 0.15$--$0.19$ is tighter because the full sample $n=2000$ is used without splitting, but it delivers no per-record attribution. HIF's row-level design therefore trades per-cell detectability for granularity; the multi-hub geometric aggregation and the confidence-floor safety minimum ($0.01$) are the design elements that dampen the long-tail cells' influence on the final score.

\subsection{Scalability and Baseline Comparisons}\label{app:scalability}
We benchmarked HIF against the \textbf{SHAP Distance} semantic fidelity metric \cite{yu2025shap}, which serves as a state-of-the-art baseline for explainability-aware evaluation. Although the SHAP Distance effectively identifies semantic drift, it suffers from severe scalability bottlenecks. HIF auditing scales linearly, $O(N \cdot K)$, where $K$ is the number of hubs. In our empirical scalability benchmark, fitting the 5-hub LSE auditor on 2,000 real records took $9.7$\,s, after which scoring 10,000 synthetic records took $1.50 \pm 0.15$\,s and scoring 100,000 records took $10.6 \pm 0.5$\,s on commodity multi-core hardware (fit amortized once; wall-clock averaged over three repeats). The marginal per-10,000-row cost at the 100,000-row scale ($10.6\,/\,10 \approx 1.1$\,s per 10k rows) confirms the linear scaling. In contrast, SHAP-based methods require $O(N \cdot D \cdot 2^D)$ operations or expensive approximations, which precludes exact evaluation at this volume. Thus, HIF provides comparable detection accuracy while maintaining the computational efficiency required to scale to high-cardinality datasets with millions of rows.
